\section{Borsuk--Ulam theorem}

\begin{definition}[Unit sphere and antipodal map]
\label{def:borsuk-ulam-antipode}
\lean{BorsukUlamTheorem.UnitSphere, BorsukUlamTheorem.antipode}
\leanok
For $n\in\mathbb N$, let $S^n$ be the unit sphere in $\mathbb R^{n+1}$. The antipodal map is $a(x)=-x$.
\end{definition}

\begin{lemma}[Continuity of the antipodal map]
\label{lem:borsuk-ulam-continuous-antipode}
\lean{BorsukUlamTheorem.continuous_antipode}
\leanok
\uses{def:borsuk-ulam-antipode}
The antipodal map on $S^n$ is continuous.
\end{lemma}

\begin{lemma}[The antipodal map is an involution]
\label{lem:borsuk-ulam-antipode-involution}
\lean{BorsukUlamTheorem.antipode_antipode}
\leanok
\uses{def:borsuk-ulam-antipode}
For every $x\in S^n$, one has $a(a(x))=x$.
\end{lemma}

\begin{theorem}[Borsuk--Ulam]
\label{thm:borsuk-ulam}
\lean{BorsukUlamTheorem.MainTheorem}
\notready
\uses{def:borsuk-ulam-antipode}
For every $n\in\mathbb N$ and every continuous map $f:S^n\to\mathbb R^n$, there exists $x\in S^n$ such that $f(x)=f(-x)$.
\end{theorem}

\begin{proof}
\uses{lem:borsuk-ulam-continuous-antipode, lem:borsuk-ulam-antipode-involution}
The Lean proof handles dimensions $0$ and $1$. In dimension $1$ it uses connectedness and the intermediate value theorem. The case $n\ge 2$ is currently left as a \texttt{sorry}.
\end{proof}
